% =====================================================================
\section{Liên hệ với Edge AI và IoT}
\label{sec:iot}
% =====================================================================

\subsection{Vì sao chi phí suy luận mới là chi phí quyết định trên thiết bị biên}
\label{sec:iot-why}

Trong học đa nhiệm, việc thêm một nhánh phụ làm tăng chi phí \emph{huấn luyện}, nhưng chi phí
\emph{suy luận} — thứ quyết định trên thiết bị triển khai — có thể hoàn toàn không đổi. Đây là
điểm quan trọng cho các ứng dụng IoT: mô hình được huấn luyện một lần trên máy chủ nhưng được
chạy liên tục trên thiết bị. Với cấu hình tại \secref{sec:arch}, mô hình suy luận của \cfg{E1}
và \cfg{E2} có \emph{đúng cùng} kiến trúc vì Rotation Head không tham gia suy luận. Vì vậy:

\begin{keybox}
Nhánh phụ không làm mô hình suy luận nhỏ hơn hay lớn hơn. Lợi ích của học đa nhiệm (nếu có)
phải đến từ \emph{chất lượng đặc trưng}, không từ việc giảm độ phức tạp kiến trúc. Mọi tuyên bố
về \enquote{tiết kiệm tài nguyên} chỉ nhờ bỏ head phụ là không chính xác về mặt kỹ thuật.
\end{keybox}

Bảng \ref{tab:inference-cost} cho thấy mức chênh lệch cụ thể: \RParamRotationHead{} tham số của
Rotation Head tương đương khoảng $0{,}018\,\%$ tổng số tham số huấn luyện — không đáng kể so
với \RParamInference{} tham số của mô hình suy luận.

\subsection{Đối chiếu với ngân sách tài nguyên của thiết bị IoT}
\label{sec:iot-budget}

Để đánh giá mức độ phù hợp với thiết bị biên, cần so số tham số đo được với ngân sách bộ nhớ
flash và RAM thực tế của các lớp thiết bị. Các hướng dẫn benchmark cho thiết bị siêu nhỏ hiện
dùng ngân sách cỡ vài trăm KiB bộ nhớ cho cả mô hình và trạng thái trung gian, trong khi các mô
hình tối ưu cho vi điều khiển được thiết kế từ đầu với ràng buộc đó~\cite{banbury2021mlperftiny,lin2019mcunet}.

\begin{table}[htbp]
  \centering
  \small
  \begin{tabular}{@{}l r l@{}}
    \toprule
    Thành phần & Dung lượng & Ghi chú \\
    \midrule
    Trọng số FP32 của mô hình suy luận & \RParamInferenceMiB{} MiB & đo được: \RParamInferenceFileMiB{} MiB \\
    Trọng số lượng tử hóa INT8 (ước lượng) & $\approx 10{,}65$ MiB & $11.169.858 \times 1$ byte \\
    Trọng số lượng tử hóa INT4 (ước lượng) & $\approx 5{,}33$ MiB & $11.169.858 \times 0{,}5$ byte \\
    \midrule
    Ngân sách flash điển hình của MCU nhỏ & vài trăm KiB & xem \cite{lin2019mcunet} \\
    \bottomrule
  \end{tabular}
  \caption{Đối chiếu dung lượng trọng số của mô hình với ngân sách bộ nhớ của thiết bị biên. Các
  dòng lượng tử hóa là \emph{ước lượng số học đơn giản} từ số tham số, chưa bao gồm chi phí
  hiệu chỉnh, activation và bộ nhớ runtime; chúng \textbf{không} phải là kết quả đã đo. Hai dòng
  cuối chỉ mang tính đối chiếu định cỡ, không phải suy luận về khả năng triển khai.}
  \label{tab:iot-budget}
\end{table}

Hệ quả định cỡ: ngay cả ở dạng lượng tử hóa INT8, mô hình này vượt xa ngân sách bộ nhớ của một
vi điều khiển nhỏ. Với các nền tảng mạnh hơn như vi điều khiển ứng dụng có vài MiB flash, hoặc
các SoC biên có bộ tăng tốc NPU, mô hình lại nằm trong dải khả thi. Việc rút gọn backbone là
điều kiện tiên quyết nếu mục tiêu là vi điều khiển lớp nhỏ; các kiến trúc hiệu quả tham số như
MobileNet cho thấy hướng đi này~\cite{howard2017mobilenets}.

\subsection{Giao thức đo độ trễ và giới hạn kết luận}
\label{sec:iot-latency}

Độ trễ đo được tại \tabref{tab:latency} có giá trị tham chiếu nội bộ, cần được đọc kèm các điều
kiện sau:

\begin{itemize}[leftmargin=1.2cm,itemsep=3pt]
  \item Phép đo chỉ tính thời gian tính toán của mô hình. Đọc ảnh, tiền xử lý và truyền dữ liệu
        được đo riêng (\RPreprocessMs{}~ms/ảnh trên CPU) và đây là thành phần thường chiếm phần
        lớn thời gian trong một pipeline thực tế.
  \item Trên GPU, p95 cao hơn trung vị đáng kể; phân bố độ trễ không đối xứng nên không nên
        dùng duy nhất giá trị trung bình để thiết kế hệ thống.
  \item Số luồng CPU được ghim cố định (\RLatThreads{} luồng) để các phép so sánh có ý nghĩa.
        Thay đổi số luồng sẽ thay đổi kết quả và cần được khai báo lại.
\end{itemize}

\begin{limbox}
\textbf{Giới hạn kết luận được nêu rõ.} Kết quả trên CPU của máy tính dùng để huấn luyện
\emph{không} phải là benchmark trên thiết bị đích. Chỉ có thể kết luận về \enquote{khả năng
triển khai thiết bị biên} khi đã đo trên chính thiết bị đó, với runtime và mức lượng tử hóa
thực tế. Báo cáo này \textbf{không} suy diễn từ số tham số hay độ trễ trên máy tính thành kết
luận về vi điều khiển.
\end{limbox}

\subsection{Hàm ý thiết kế cho ứng dụng IoT}
\label{sec:iot-implications}

Nếu hiệu ứng tích cực của nhiệm vụ phụ được xác nhận ở quy mô lớn hơn, hàm ý thực tiễn cho
ứng dụng IoT là: \emph{chi phí một lần khi huấn luyện để đổi lấy chất lượng suy luận tốt hơn ở
cùng kích thước mô hình}. Với các hệ thống thu thập nhãn đắt đỏ (ví dụ cảm biến công nghiệp
cần chuyên gia gán nhãn), việc tận dụng tín hiệu tự giám sát từ dữ liệu chưa gán nhãn có thể
rẻ hơn đáng kể so với việc gán nhãn thêm. Ngược lại, nếu hiệu ứng không được xác nhận, kết luận
hữu ích cũng vẫn có giá trị: nó cho thấy trong điều kiện cụ thể đã khảo sát, chi phí huấn luyện
tăng thêm \emph{không} đổi lấy chất lượng suy luận tốt hơn.

Một điểm cần nhấn mạnh cho thiết kế hệ thống: \textbf{không có} thay đổi nào về đường suy luận
giữa \cfg{E1} và \cfg{E2}. Điều này có nghĩa là bất kỳ hạ tầng triển khai nào đã chấp nhận
\cfg{E1} đều chấp nhận \cfg{E2} mà không cần thay đổi runtime, bộ nhớ hay giao thức. Rủi ro
triển khai của việc áp dụng học đa nhiệm trong trường hợp này là rất thấp.
